\documentclass{beamer}

\usepackage[utf8]{inputenc}
\usepackage{amssymb}
\usepackage{amsfonts}

\title{Collaborative Folding Schemes \\ \emph{\small A proposal}}
\author{Daniel Barrero}
\institute{Universiad de los Andes}


\begin{document}
\frame{\titlepage}

\begin{frame}{Arguments of Knowledge}
$\pvr$ wants to convince $\vfr$ that it knows a witness $w$ to a statement
$x$.
\begin{example}
	\begin{enumerate}
		\item
		\[
			\left\{
			\begin{array}{l}
				x= \mbox{\emph{``There exists
				$\sigma$ such that
				SHA256$(\sigma)$
				is equal to 7''.}}\\
				w=\  \mbox{Some string such that
				SHA256$(w) = 7$.} 
			\end{array}
			\right.
		\]
		\item
		\[
			\left\{
				\begin{array}{l}
					x= \mbox{\emph{``The dodecahedron
					defines a hamiltonian graph''.}}\\
					w= \mbox{A hamiltonian cycle in
					the dodecahedron.}
				\end{array}
			\right.
		\]
	\end{enumerate}
\end{example}
\end{frame}

%nextframe
\begin{frame}{Arguments of Knowledge}
\begin{example}
	\begin{enumerate}
		\setcounter{enumi}{2}
		\item
		\[
			\left\{
				\begin{array}{l}
					x= \mbox{``There exists more than
					one $Q$-short vector}\\
					\mbox{in the lattice
					$Bh+t$''.}\\
					w= \mbox{A HAWK signature.}
				\end{array}
			\right.
		\]
	\end{enumerate}
\end{example}

\begin{block}{Proofs of knowledge}
	A \emph{proof of knowledge} is some message $\pi$ that $\pvr$ sends
	to $\vfr$ to convince her that he knows $w$. 
	\begin{itemize}
		\item The \textbf{trivial proof} Is when $\pi= w$. It is
			impractical because (1) It could occupy too much
			space, and (2) Assumes $\vfr$ has the same
			computational power as $\pvr$.
	\end{itemize}
\end{block}
\end{frame}

%nextframe
\begin{frame}{Arguments of Knowledge}
	\Large
	\[
		\tikzfig{tradark}
	\]
	\normalsize
	\begin{definition}
		An \emph{argument of knowledge (ARK)} between a prover $\pvr$
		and a verifier $\vfr$ is an interaction where
		\begin{itemize}
			\item $\pvr$ has inputs $x$ and $w$, $\vfr$ has input $x$
			\item $\pvr$ outputs $\pi$ and sends it to
				$\vfr$
			\item $\vfr$ outputs $0/1$ dependending on whether
				it rejects or accepts $\pi$
		\end{itemize}
	\end{definition}
\end{frame}

%nextframe
\begin{frame}{ARKs - Formal properties}
	In cryptography, the statement-witness pairs $(x,w)$ must take values
	in an NP relation $\kre$. Also, the protocol must be
	\emph{complete and sound}. Informally, this means
	\[
		(x,w) \in \kre \iff \mbox{$\vfr$ accepts $\pi$}
	\]
	where
	\begin{itemize}
		\item \textbf{Completeness} corresponds to
			\[
				(x,w) \in \kre \Rightarrow
				\mbox{$\vfr$ accepts $\pi$}
			\]
		\item \textbf{Soundness} ``corresponds'' to
			\[
				\mbox{$\vfr$ accepts $\pi$} \Rightarrow
				(x,w) \in \kre
			\]
	\end{itemize}
	\begin{quote}
		WARNING!! \emph{Soundness is actually defined in terms of a malicious
		prover $\pvr^*$ and an aditional algorithm $\mathcal{E}$
		called the \textbf{extractor}}
	\end{quote}
\end{frame}

%nextframe
\begin{frame}{Reductions of Knowledge}
	Contemporary arguments of knowledge such as \textrm{LatticeFold}
	require a generalization of ARKs called \emph{reductions of knowledge
	(RoK)}.
	\begin{definition}
		Let $\kre_1, \kre_2$ be NP relations of statement-witness
		pairs. A \emph{reduction of knowledge} from $\kre_1$ to
		$\kre_2$ is a protocol $\prod_{\pvr,\vfr}$ such that
		\begin{itemize}
			\item $\pvr$ has input $(x_1,w_1) \in \kre_1$
			\item $\vfr$ has input $x_1$
			\item They interact and as a result
				\[
					\left\{
						\begin{array}{l}
							\mbox{$\pvr$ and $\vfr$
							both output $x_2$}\\
							\mbox{$\pvr$ outputs
							$w_2$}\\
							(x_2, w_2) \in \kre_2.
						\end{array}
					\right.
				\]
			\item $\prod_{\pvr,\vfr}$ is complete and *sound*
		\end{itemize}
	\end{definition}
\end{frame}

%nextframe
\begin{frame}{Reductions of Knowledge}
	\large
	\ctikzfig{rok}
\end{frame}
\end{document}
